\documentclass[10pt]{article}
\usepackage[T1]{fontenc}
\usepackage[utf8]{inputenc}
\usepackage[french]{babel}
\usepackage{amsmath,amssymb,mathrsfs,mathpazo}
\usepackage{geometry}
\geometry{a4paper,margin=2cm}
\usepackage{graphicx,array,tabularx,tabularray,colortbl,multirow,multicol,booktabs,enumitem,pifont,ulem,fancybox,diagbox,cancel}
\usepackage{pstricks,pstricks-add,pst-plot,pst-tree,pst-node,pst-text,pst-func,pst-math,pst-bspline,pst-3dplot,pst-eucl,pst-solides3d}
\let\pscalculate\undefined
\usepackage{pst-bezier}
\usepackage{tikz,pgfplots,tkz-tab}
\usetikzlibrary{arrows,arrows.meta,patterns,calc,positioning,trees,decorations.markings,shapes.geometric,shapes.misc,shapes.arrows,intersections,decorations.pathmorphing,decorations.pathreplacing,backgrounds,shadows,babel}
\pgfplotsset{compat=1.18}
\usepgfplotslibrary{fillbetween}
\usepackage[np]{numprint}
\usepackage{listings,makecell,pas-tableur,twoopt}
\let\barre\undefined
\usepackage[pstricks]{tablvar}
\definecolor{darkgreen}{RGB}{0,100,0}
\providecommand{\red}{\color{red}}
\providecommand{\blue}{\color{blue}}
\providecommand{\green}{\color{green}}
\providecommand{\black}{\color{black}}
\providecommand{\white}{\color{white}}
\providecommand{\gray}{\color{gray}}
\providecommand{\cyan}{\color{cyan}}
\providecommand{\magenta}{\color{magenta}}
\providecommand{\yellow}{\color{yellow}}
\providecommand{\violet}{\color{violet}}
\providecommand{\orange}{\color{orange}}
\providecommand{\euro}{\texteuro}
\providecommand{\EUR}{\texteuro}
\providecommand{\ieme}{\textsuperscript{e}}
\providecommand{\degres}{\ensuremath{^\circ}}
\providecommand{\degree}{\ensuremath{^\circ}}
\providecommand{\vv}[1]{\overrightarrow{#1}}
\providecommand{\Coord}[1]{#1}
\providecommand{\decosix}{}
\providecommand{\starredbullet}{\ensuremath{\bullet}}
\pagestyle{empty}
\setlength{\parindent}{0pt}
\usepackage{adjustbox}\DeclareUnicodeCharacter{2236}{:}\DeclareUnicodeCharacter{2190}{\ensuremath{\leftarrow}}
\providecommand{\R}{\mathbb{R}}
\providecommand{\N}{\mathbb{N}}
\providecommand{\Z}{\mathbb{Z}}
\providecommand{\Q}{\mathbb{Q}}
\providecommand{\D}{\mathbb{D}}
\providecommand{\C}{\mathbb{C}}
\providecommand{\e}{{\mathrm{e}}}
\providecommand{\vect}[1]{\overrightarrow{#1}}
\providecommand{\vectt}[1]{\overrightarrow{\mathrm{#1}}}
\providecommand{\vvt}[1]{\overrightarrow{\mathrm{#1}}}
\providecommand{\np}[2][]{{#2~#1}}
\providecommand{\vv}[1]{{\overrightarrow{#1}}}
\renewcommand{\d}{{\mathrm{d}}}
\providecommand{\degres}{{^\circ}}
\providecommand{\degre}{{^\circ}}
\providecommand{\blue}{\color{blue}}
\providecommand{\red}{\color{red}}
\providecommand{\enskip}{\quad}
\providecommand{\llbracket}{\lbrack\!\lbrack}
\providecommand{\rrbracket}{\rbrack\!\rbrack}
\providecommand{\euro}{{\text{€}}}
\providecommand{\no}{{\text{n°}}}
\providecommand{\interval}[2]{{[#1;#2]}}
\providecommand{\intervFO}[2]{{[#1;#2[}}
\providecommand{\dotfill}{\ldots}

\providecommand{\ds}{\displaystyle}
\providecommand{\textding}[1]{#1}
\providecommand{\Oij}{\(\mathrm{O};\vec{\imath},\vec{\jmath}\)}
\providecommand{\Oijk}{\(\mathrm{O};\vec{\imath},\vec{\jmath},\vec{k}\)}
\providecommand{\Ouv}{\(\mathrm{O};\vec{u},\vec{v}\)}
\providecommand{\intervFF}[2]{\ensuremath{[#1;#2]}}
\renewcommand{\euro}{\ifmmode\text{\texteuro}\else\texteuro\fi}
\DeclareUnicodeCharacter{00B0}{\ifmmode^{\circ}\else\textdegree\fi}
\DeclareUnicodeCharacter{20AC}{\ifmmode\text{\texteuro}\else\texteuro\fi}
% Énoncés : Éducation nationale. Transcription et corrigés : APMEP, Denis Vergès et contributeurs. Découpage et classement : CoopMaths, Sébastien Lozano, Éric Elter et contributeurs. Adaptation de la présentation : Maths Explorer.
% CC BY-SA : https://coopmaths.fr/www/about/
% https://coopmaths.fr/alea/static/bac/2023/tex/bac_2023_09_sujet1_polynesie_2_cor.tex
\begin{document}
\noindent\textit{APMEP / CoopMaths. CC BY-SA. Présentation adaptée par Maths Explorer.}\par\medskip
\medskip

\textbf{Partie A}

\medskip

On considère la fonction $f$ définie sur $\R$ par
$f(x) = \left(x + \dfrac12\right)\e^{-x} + x.$



\begin{enumerate}
\item On détermine les limites de $f$ en $- \infty$ et en $+ \infty$.

\begin{list}{\textbullet}{}
\item Limite en $-\infty$

$\left .
\begin{array}{r}
\left .
\begin{array}{l}
\ds\lim_{x\to -\infty} \left (x+\dfrac{1}{2}\right )=-\infty\\
\ds\lim_{x\to -\infty} \e^{-x} = +\infty\\
\end{array}
\right \}
\ds\lim_{x\to -\infty} \left (x+\dfrac{1}{2}\right ) \e^{-x}=-\infty\\
\ds\lim_{x\to -\infty} x = -\infty
\end{array}
\right \}
\text{ donc }\ds\lim_{x\to -\infty} f(x)=-\infty$

\item Limite en $+\infty$

$f(x)= x\e^{-x} + \dfrac{1}{2}\e^{-x} +x$.
On sait que $\ds\lim_{x\to +\infty} \e^{-x}=0$ et que\\

$\ds\lim_{x\to +\infty}  \dfrac{\e^{x}}{x}=+\infty$ donc $\ds\lim_{x\to +\infty}  \dfrac{x}{\e^{x}}=0$ donc $\ds\lim_{x\to +\infty}  x\e^{-x}=0$.

$\left .
\begin{array}{l}
\ds\lim_{x\to +\infty} x \e^{-x}=0\\
\ds\lim_{x\to +\infty} \dfrac{1}{2}\e^{-x} = 0\\
\ds\lim_{x\to +\infty} x = +\infty
\end{array}
\right \}
\text{ donc }\ds\lim_{x\to +\infty} f(x)=+\infty$
\end{list}

\item On admet que $f$ est deux fois dérivable sur $\R$.
	\begin{enumerate}
		\item 
Pour tout $x$ de $\R$, $f(x)= \left (x+\dfrac{1}{2}\right )\e^{-x}  +x$ donc

$f'(x)=1\times \e^{-x} + \left (x+\dfrac{1}{2}\right )\times (-1) \e^{-x} + 1
= \left ( -x +\dfrac{1}{2}\right )\e^{-x}+1$, et donc

$f''(x)= -1\times \e^{-x} + \left (-x+\dfrac{1}{2}\right )\times (-1) \e^{-x}
= \left ( x-\dfrac{3}{2}\right )\e^{-x}$. 

		\item 
Le signe de $f''$ donne les variations de $f'$.

Pour tout $x$, $\e^{-x}>0$ donc $f''(x)$ est du 	signe de $\left (x-\dfrac{3}{2}\right )$.

\begin{list}{\textbullet}{}
\item Si $x<\dfrac{3}{2}$, $f''(x)<0$ donc $f'$ est décroissante;
\item Si $x>\dfrac{3}{2}$, $f''(x)>0$ donc $f'$ est croissante;
\item Si $x=\dfrac{3}{2}$, $f''(x)=0$ donc $f'$ admet un minimum égal à \\
$f'\left (\dfrac{3}{2}\right ) = \left (-\dfrac{3}{2} + \dfrac{1}{2}\right )\e^{-\frac{3}{2}}+1=1-\e^{-\frac{3}{2}}$.
\end{list}

		\item 
La fonction $f'$ admet pour minimum

$f'\left (\dfrac{3}{2}\right ) =1-\e^{-\frac{3}{2}} \approx 0,78>0$; donc   pour tout $x \in \R,\: f'(x) > 0$.
		
		\item 
\begin{list}{\textbullet}{}
\item La fonction $f$ est dérivable sur $\R$ donc continue sur $\R$.
\item Pour tout réel $x$, $f'(x)>0$ donc la fonction $f$ est strictement croissante sur $\R$.
\item $\ds\lim_{x\to -\infty} f(x)=-\infty$
\item $\ds\lim_{x\to +\infty} f(x)=+\infty$
\end{list}		

Donc, d'après le corollaire du théorème des valeurs intermédiaires, l'équation $f(x) = 0$ admet une solution unique dans $\R$.
		
		\item 
On appelle $\alpha$ la solution de l'équation $f(x) = 0$.

$\left .
\begin{array}{l}
f(-1)\approx -2,36<0\\
f(0)=0,5>0
\end{array}
\right \}
\text{ donc } \alpha\in \left [-1\;;\;0\strut\right ]$
		
$\left .
\begin{array}{l}
f(-0,3)\approx -0,03<0\\
f(-0,2)\approx 0,17>0
\end{array}
\right \}
\text{ donc } \alpha\in \left [-0,3\;;\;-0,2\strut\right ]$

$\left .
\begin{array}{l}
f(-0,29)\approx -0,009<0\\
f(-0,28)\approx 0,011>0
\end{array}
\right \}
\text{ donc } \alpha\in \left [-0,29\;;\;-0,28\strut\right ]$
		
$\left .
\begin{array}{l}
f(-0,286)\approx -0,001<0\\
f(-0,285)\approx \np{0,0009}>0
\end{array}
\right \}
\text{ donc } \alpha\in \left [-0,286\;;\;-0,285\strut\right ]$		

Donc $-0,285$ est une valeur approchée à $10^{-3}$ de la solution de l'équation $f(x)=0$.
		
	\end{enumerate}
\end{enumerate}

\medskip
\textbf{Partie B}
\medskip

On considère une fonction $h$, définie et dérivable sur $\R$, ayant une expression de la forme 

\hfill$h(x)= (ax + b)\e^{-x} + x$, \:où $a$ et $b$\: sont deux réels.\hfill\null

\begin{list}{\textbullet}{Dans un repère orthonormé ci-après figurent:}
\item la courbe représentative $\mathcal{C}_h$ de la fonction $h$ ;
\item les points A et B de coordonnées respectives $(-2~;~-2,5)$ et $(2~;~3,5)$.
\end{list}

\begin{center}
\psset{unit=0.8cm}
\begin{pspicture*}(-5.9,-4.9)(5.9,5.9)
\psaxes[linewidth=1.25pt,labelFontSize=\scriptstyle](0,0)(-5.9,-4.9)(5.9,5.9)
\psplot[plotpoints=500,linewidth=1.25pt]{-4}{4}{1.5 x mul 0.5 add}
\psplot[plotpoints=2000,linewidth=1.25pt,linecolor=cyan]{-4}{5}{x 0.5 add 2.71828 x exp div x add}
\psdots(-2,-2.5)(2,3.5)
\uput[r](-2,-2.5){A}\uput[r](2,3.5){B}\uput[dr](3.8,4){\cyan $\mathcal{C}_h$}
\psplot[plotpoints=2000,linewidth=1.25pt,linecolor=red,linestyle=dashed]{-4}{5}{x 0.816 mul 0.736 add}
\psline[linecolor=red,linestyle=dashed,showpoints](1.5,0)(1.5,1.95)
\uput[d](1.5,0){\red $1,5$}
\end{pspicture*}
\end{center}

\begin{enumerate}
\item Avec la précision permise par le graphique, on peut conjecturer que la courbe représentative de la fonction $h$ admet un unique point d'inflexion d'abscisse  $1,5$.

\item On admet que $h''(x) = - \dfrac32 \e^{- x} + x\e^{- x}$.

$h''(x)=0
\iff
-\dfrac{3}{2}\e^{-x}+x\e^{-x}=0
\iff
\left (x-\dfrac{3}{2}\right )\e^{-x}=0
\iff
x-\dfrac{3}{2}=0
\iff
x=1,5$

La conjecture est donc vérifiée.

\item Déterminer une équation de la droite (AB) c'est chercher l'ensemble des points de coordonnées $(x~;~y)$ tels que:
$\dfrac{y-y_{\text A}}{x-x_{\text A}} = \dfrac{y_{\text B}-y_{\text A}}{x_{\text B}-x_{\text A}}$.

$\begin{gathered}\dfrac{y-y_{\text A}}{x-x_{\text A}} = \dfrac{y_{\text B}-y_{\text A}}{x_{\text B}-x_{\text A}}
\iff
\dfrac{y+2,5}{x+2} = \dfrac{3,5 + 2,5}{2 + 2}
\iff
\dfrac{y+2,5}{x+2} = \dfrac{6}{4}
\iff
y+2,5 = \dfrac{6}{4}\left (x+2\right )\\
\phantom{\dfrac{y-y_{\text A}}{x-x_{\text A}} = \dfrac{y_{\text B}-y_{\text A}}{x_{\text B}-x_{\text A}}}
\iff
y=\dfrac{3}{2}x +3-2,5
\iff
y=1,5x + 0,5\end{gathered}$

La droite (AB) a pour équation: $y =1,5x + 0,5$.

\item La droite (AB) est tangente à la courbe représentative de la fonction $h$ au point d'abscisse 0, donc $h'(0)=1,5$.

$h'(x)=a\times \e^{-x} +  (ax+b)\times (-1)\e^{-x}+1
=  (-ax+a-b )\e^{-x}+1$

$h'(0)=1,5
\iff
(a-b)\e^{0}+1=1,5
\iff
a-b=0,5$

La droite (AB) coupe l'axe des ordonnées au point d'ordonnée $0,5$ donc $h(0)=0,5$.

$h(0)=0,5
\iff
(0+b)\e^{0}+0=0,5
\iff
b=0,5$

Comme $a-b=0,5$, on en déduit que $a=1$.

Donc $h(x)=(x+0,5)\e^{-x}+x=f(x)$.
\end{enumerate}

\bigskip

\end{document}
